\documentclass{article}

\textwidth=6.5in
\textheight=8.5in
\voffset=-54pt
\oddsidemargin=0in
\evensidemargin=0in
\setlength{\parskip}{8pt}
\setlength{\parindent}{0pt}

\usepackage{amsmath, amssymb}
\usepackage{ifthen}

\usepackage[inline]{enumitem}
\setlist{topsep=0pt, parsep=8pt}

\usepackage[rgb,table]{xcolor}\convertcolorsUfalse
\usepackage{graphicx}

% change hyperref options
\PassOptionsToPackage{hyphens}{url}
\usepackage[bookmarks,colorlinks,linkcolor=black,citecolor=blue,pdfstartview=FitH,urlcolor=blue]{hyperref}
\hypersetup{pdfpagemode=UseNone}
\usepackage{bookmark}

\usepackage{graphicx}
\usepackage{multicol}
\usepackage{framed}
\usepackage{array}

\usepackage{icomma}

\usepackage{soul}

\usepackage{pgfplots}
\pgfplotsset{compat=1.18}  
\pgfplotscreateplotcyclelist{mycolor}{black, blue, red, green!70!black, magenta, purple, cyan, orange }  
\pgfplotsset{
  disabledatascaling,
  cycle list=tikzcolor, 
  axis lines=middle,
  every axis plot/.style={no marks, thick},
  every axis/.style={
    enlarge x limits=false,
    enlarge y limits=auto,
    xlabel style={at={(current axis.right of origin)},anchor=west},
    ylabel style={at={(current axis.above origin)},anchor=south},
    % extra x and y ticks for asymptotes
    extra x tick style={grid=major, grid style={dashed,black}},
    extra y tick style={grid=major, grid style={dashed,black}},
    /pgf/number format/1000 sep={},
  },    
  tick label style = {font=\footnotesize, fill=\currentbgcolor, inner sep=0pt},    
  xticklabel shift=1pt,    
  scaled ticks=false,
  scaled y ticks=false,
  unbounded coords=jump,
  samples=101,
  minor tick length=0,
  scatter/classes={
    open={mark=*,draw=black,fill=white, mark size=2, thin},
    closed={mark=*,draw=black,fill=black, mark size=2, thin}
  },
  width=3in,
}
\pgfplotsset{open/.style={only marks, mark=*, fill=white, thin, mark size=2}}
\pgfplotsset{closed/.style={only marks, mark=*, draw=black, fill=black,
    thin, mark size=2}}
\pgfplotsset{labeled/.style={nodes near coords, only marks, mark=*, draw=black, fill=black, thin, mark size=2, point meta=explicit symbolic}}

\pgfplotsset{gridgraph/.style={clip=false, axis line style={shorten
      >=-8pt}, xlabel style={xshift=8pt}, ylabel style={yshift=8pt},
    enlargelimits=false, grid=major, tick style={black}}} 

\usepackage{tikz}
\usetikzlibrary{
  matrix,%
  % enables the snake paths
  decorations.pathmorphing,%
  % enables bigger arrows
  decorations.markings,%
  decorations.pathreplacing,%
  decorations.text,%
  calc,%
  patterns,%
  shapes.geometric,%
  arrows,
  decorations.shapes,
  positioning,
  plotmarks,
  intersections
}
\tikzset{>=latex} % sets default arrow type in tikz
\tikzset{font=\normalsize}
\tikzset{mark size=1.5pt, mark options=thin}
\tikzset{pin distance=4pt,
  every pin edge/.style={<-, thin, shorten <= -2pt}}

\interfootnotelinepenalty=10000
\renewcommand{\thefootnote}{(\arabic{footnote})}

\usepackage{multirow, bigdelim}

% change to \usepackage[sols]{handout} to see solutions
\usepackage[sols]{handout}

\newcommand\iscurrentenumi[1]{%
  \ifthenelse{\equal{#1}{\theenumi}}{}{#1}}
\setenumerate[1]{ref=\#\arabic*}
\setenumerate[2]{ref=\protect\iscurrentenumi{\theenumi}(\alph*)}
\newcommand\enumiiprefix[2]{%
  \ifthenelse{{\equal{#1}{\theenumi}}\AND{\equal{#2}{(\alph{enumii})}}}{}{}%
  \ifthenelse{{\equal{#1}{\theenumi}}\AND{\NOT{\equal{#2}{(\alph{enumii})}}}}{#2}{}%
  \ifthenelse{\equal{#1}{\theenumi}}{}{#1#2}%
}
\setenumerate[3]{ref=\protect\enumiiprefix{\theenumi}{(\alph{enumii})}\roman*}

\newlength{\tipmargin}
\makeatletter

% underline links               
\let\@href=\href
\renewcommand{\href}[2]{\@href{#1}{\ul{#2}}}

\renewenvironment{framed}{
  \setlength{\tipmargin}{\textwidth-\linewidth}
  \advance\@totalleftmargin-\tipmargin
  \def\FrameCommand{\hspace{\tipmargin}\fbox}%
  \advance\linewidth\tipmargin
  \MakeFramed {\advance\hsize-\width \FrameRestore}%
}
{\endMakeFramed}

\makeatother

\definecolor{purple}{rgb}{0.5,0,1.0}

\newcounter{imagecounter}
\newcounter{columncounter}
\newenvironment{imagetable}[3][]{%
  \setcounter{imagecounter}{0}
  \setcounter{columncounter}{#2}
  \addtocounter{columncounter}{-1}
  \newcolumntype{i}{>{\raisebox{-1ex-\height}{\makebox[1em]{\hfill\stepcounter{imagecounter}(#3{imagecounter})}}\hspace{4pt}\begin{minipage}[t]{\linewidth/#2-1em-\tabcolsep*3}\arraybackslash\vspace{0pt}}l<{\end{minipage}}}
\begin{tabular}{i*{\value{columncounter}}{#1i}}}
{\end{tabular}}  

\usepackage{fancyhdr}
\lhead{\textcolor{white}{This content is protected and may not be shared, uploaded, or distributed.}}
\rhead{}
\rfoot{\vspace*{\fill}\footnotesize\textcopyright{} \emph{Harvard Math Ma}}
\renewcommand{\headrulewidth}{0pt}
\pagestyle{fancy}
\begin{document}

\htitle{More on Functions, Average Rates of Change}

\begin{problemonly}
  \begin{framed}

You should be able to:
    \begin{itemize}[topsep=0pt,partopsep=0pt,parsep=8pt,itemsep=0pt]

    \item Write and interpret expressions in function notation.

    \item Compute an average rate of change, interpret it in the context of a situation, and represent it graphically.

    \item Explain the difference between average speed, average velocity, instantaneous speed, and instantaneous velocity.

    \item Explain how concavity can be characterized in terms of secant line segments.

    \end{itemize}
  \end{framed}
\end{problemonly}

\begin{enumerate}
\item 

  \note{The point here is just to have students say what average speed means before they start thinking about instantaneous speeds. That way, if they later try to calculate an average speed by averaging some instantaneous speed, we can come back to this to remember that isn't the definition of average speed.}

  \begin{problem}[0.5in]
    \emph{Warm up.} Your friend is preparing for a marathon and goes
    for a training run. What information would you need to be able to
    calculate your friend's average speed on this run?
  \end{problem}

  \begin{solution}
    You need to know the distance your friend ran and the amount of
    time it took. Then, average speed =
    $\dfrac{\text{distance}}{\text{elapsed time}}$. 
  \end{solution}

\item  \label{gottlieb-1.2.12-followup}

  \note{Robin explains: ``This will make clear the connection between
    words and functional notation, input and output, and modeling.
    This problem grew out of the observation that many times students
    were memorizing the definition of derivative and didn't see much
    harm in the `minor' mistake of writing $f(x) + h$ when they should
    have written $f(x + h)$. This problem is getting at the fact that
    they MEAN totally different things.''}

  In \href{https://people.math.harvard.edu/~jjchen/mathma/docs/homework/PS02.html}{Problem Set 2}, {{\#6}}, you looked at a road trip east on Route 66
  from Flagstaff, Arizona through New Mexico into Oklahoma. The road
  trippers entered Gallup, New Mexico $T$ hours into the trip.

  Let $f(t)$ be the number of miles they travel in the first $t$ hours of their trip, and let $g(t)$ be their speed $t$ hours into the trip.

  \begin{enumerate}
  \item
    \begin{problem}[\stretch{1.5}]
      How are these two functions related to instruments on the car's dashboard?
    \end{problem}

    \begin{solution}
      $f(t)$ can be thought of as the change in the car's odometer reading since the start of the trip, and $g(t)$ is the speedometer reading $t$ hours into the trip.
    \end{solution}

  \item Express the following quantities in terms of $f$ or $g$, whichever
    is appropriate. 

    \begin{enumerate}
    \item \label{gallup-instantaneous-speed-example}

      \note{If students are unsure, I will have them think about a concrete example of $T$, like $T = 5$ (and I will highlight how helpful this strategy is).}

      \begin{problem}[\vfill]
        The car's speed exactly 2 hours before entering Gallup.
      \end{problem}

      \begin{solution}
        \boxed{$g(T - 2)$}
      \end{solution}

    \item

      \note{ Students may want to say something like $\dfrac{g(2) + g(5)}{2}$, so you should be ready to talk about why this isn't sensible. For example, what if the car was paused at $t = 2$ and $t = 5$ but driving quickly in the middle?

        I will try to remember to mention the difference between average and instantaneous speed here (to contrast what was being asked in \ref{gallup-instantaneous-speed-example} with what's being asked in this part.}

      \begin{problem}[\vfill]
        The car's average speed between hour 2 and hour 5 of the trip.
      \end{problem}

      \begin{solution}
        \boxed{\frac{f(5) - f(2)}{5-2}}
      \end{solution}

    \item
      \begin{problem}[\vfill]
        The car's average speed in the half hour after getting to Gallup, New Mexico.
      \end{problem}

      \begin{solution} 
        \boxed{\frac{ f(T+0.5) -f(T)}{T+0.5-T}}
      \end{solution} 

    \end{enumerate}

  \item
    \begin{problem}[\vfill\newpage]
      In this situation, what do $g(T + 2)$ and $g(T) + 2$ represent? (Are they the same thing?)
    \end{problem}

    \begin{solution}
      The expression $g(T+2)$ represents the car's speed 2 hours after getting to Gallup.
      This is different from $g(T)+2$, which represents 
      the speed which is 2 miles per hour faster than the car's speed when it reached Gallup. Observe
      that the 2s have different units in the different expressions. 
      In $g(T+2)$, 2 has units of hours.
      In $g(T)+2$, 2 has units of miles per hour.
    \end{solution}

  \item 

    \begin{problem}[\vfill]
      Interpret $\displaystyle \frac{ f(T +5) - f(T)}{5}$ in words. (What does it express about the road trip?)
    \end{problem}

    \begin{solution} 
      The expression is a quotient and so we should be thinking ``is this describing a rate?'' It is helpful to consider the different pieces represent. $f(T+5)$ is the total distance traveled when they are $5+T$ hours into their trip and $f(T)$ is the distance traveled when they are $T$ hours into their trip. Here the rate represents \boxed{ \text{the average speed between hour $T$ and hour $T+5$.}}
    \end{solution}

  \end{enumerate}

\item  \label{mass-population}

  \note{This problem has two goals:
    \begin{itemize}[nosep]
    \item To see an average rate of change in a different context.
    \item To establish the key idea that an average rate of change can
      be visualized as slope of a secant line.
    \end{itemize}}

  Based on \href{http://factfinder.census.gov/servlet/GCTTable?_bm=y&-geo_id=01000US&-ds_name=PEP_2008_EST&-_lang=en&-redoLog=false&-mt_name=PEP_2008_EST_GCTT1_US40&-format=US-40&-CONTEXT=gct}{U.S. Census data}, the population of Massachusetts between 2000 and 2008 can be modeled by the function $P(t) = 6360 + 60 t - 15t^2 + 1.2 t^3$, where $P(t)$ gives the number of thousands of people in Massachusetts $t$ years after July 1, 2000.

  \begin{center}
    \begin{tikzpicture}
      \begin{axis}[width=4.5in, height=3in, xmax=8.2, ymin=6325, axis
        y discontinuity=parallel, ytick={6360, 6400, ..., 6600}]
        \addplot+[domain=0:8] {6360+60*x-15*x^2+1.2*x^3}; 
        \solonly{
          \addplot[closed] coordinates {(1,6406.2) (3,6437.4)};
          \draw[thick, ->] (1,6406.2) -- (3, 6406.2);
          \draw[thick, ->] (3,6406.2) -- (3, 6437.4);
          \draw[decorate, decoration={brace, amplitude=5pt, mirror}, xshift=0.1cm]
            (3,6406.2) -- node[right=0.1cm] {$\Delta P = P(3)-P(1)$} (3, 6437.4);
          \draw[decorate, decoration={brace, amplitude=5pt, mirror}, yshift=-0.1cm]
            (1,6406.2) -- node[below=0.1cm] {$\Delta t$} (3, 6406.2);
          \draw[thick] (1,6406.2) -- (3, 6437.4);
        }
      \end{axis}
    \end{tikzpicture}
  \end{center}

  \begin{enumerate}
  \item

    \note{Students may be tempted to say 6360 rather than 6360
      thousand.} 

    \begin{problem}[\stretch{0.5}]
      According to this model, what was the population of Massachusetts on
      July 1, 2000? 
    \end{problem}

    \begin{solution}
      Since $P(0) = 6360$, the population of Massachusetts on July 1, 2000
      was $6360$ thousand people, or \boxed{6.36 \text{ million people}}
    \end{solution}

  \end{enumerate}

  \probonly{Copy the graph onto the board, and draw in your answers to
    the next few parts.}

  \begin{enumerate}[resume]
  \item

    \note{If I remember, I will denote this $\Delta P$.}

    \begin{problem}[\vfill]
      Write an expression giving the change in population from July 1, 2001 to July 1, 2003. How could you visualize this quantity on the graph above?
    \end{problem}

    \begin{solution}
      The population on July 1, 2001 is $P(1)$, and the population on July 1, 2003 is $P(3)$. 
      The change in population between these two dates is $\boxed{P(3)-P(1).}$ Visually, this is 
      represented by a vertical distance on the graph.
    \end{solution}

  \item

    \note{And then this is $\dfrac{\Delta P}{\Delta t}$.}

    \begin{problem}[\vfill\newpage]
      Write an expression giving the average rate of change of the population from July 1, 2001 to July 1, 2003 (with respect to time). What are the units of your answer? How could you visualize this quantity on the graph above?
    \end{problem}

    \begin{solution}
      The average rate of change of the population is 
      $\dfrac{\text{change in population}}{\text{elapsed time}}$ or $\dfrac{\Delta P}{\Delta t}$. 
      This means that the average rate of change between July 1, 2001 and July 1, 2003 is
      $\boxed{\frac{P(3)-P(1)}{3-1},}$ with units of thousands of people per year. Visually, this
      is represented by the slope of a secant line segment on the graph.
    \end{solution}

  \item

    \note{It's much easier to do this graphically than it is to
      calculate all 3 rates!}

    \begin{problem}[\stretch{1.5}]
      Put the following quantities in increasing order. (All 3 rates
      are with respect to time.) 
      \begin{enumerate}[label=\Alph*., itemsep=0pt, parsep=0pt]
      \item The average rate of change of the population between July 1, 2001 and July 1, 2002.
      \item The average rate of change of the population between July 1, 2001 and July 1, 2003.
      \item The average rate of change of the population between July 1, 2005 and July 1, 2006.
      \end{enumerate}
    \end{problem}

    \begin{solution}
      Graphically, these quantities, represent the slopes of 3 different secant line segments.
      Looking at the slopes, we can say that
      $\boxed{C < B < A.}$

      \begin{center} 
      \begin{tikzpicture}
        \begin{axis}[width=4.5in, height=3in, xmax=8.2, ymin=6325, axis
          y discontinuity=parallel, ytick={6360, 6400, ..., 6600}]
          \addplot+[domain=0:8] {6360+60*x-15*x^2+1.2*x^3}; 
          \draw[red, very thick] (1,6406.2) -- (2, 6429.6);         
          \draw[red, very thick] (1,6406.2) -- (3, 6437.4);
          \draw[red, very thick] (5,6435) -- (6, 6439.2);
        \end{axis}
      \end{tikzpicture}
      \end{center}
    \end{solution}

  \end{enumerate}

\item  \label{concavity-secant-line-characterization}

  \begin{problem}[\vfill]
    So far, we've described concavity only qualitatively. We can give a more precise definition using secant lines. Draw a few examples of functions that are concave up on an interval $[a, b]$, as well as a few examples of functions that are concave down on an interval $[a, b]$.
  \end{problem}

  \begin{solution}
    \begin{center}
      \begin{tikzpicture}
        \begin{axis}[hide x axis, hide y axis, clip=false, width=3in, height=2in]
          \addplot+[domain=-3:3] {x^2}; %
          \pgfplotsforeachungrouped \a in {-2.1, 0, 1.75} { %
            \pgfmathsetmacro{\h}{0.75}%
            \pgfmathsetmacro{\x}{\a-\h}%
            \pgfmathsetmacro{\xx}{\a+\h}%
            \pgfmathsetmacro{\y}{\x*\x}%
            \pgfmathsetmacro{\yy}{\xx*\xx}%
            \edef\temp{\noexpand\draw[blue, thick] (\x, \y) -- (\xx, \yy);}%
            \temp%
          }
        \end{axis}
      \end{tikzpicture}
      \begin{tikzpicture}
        \begin{axis}[hide x axis, hide y axis, clip=false, width=3in, height=2in]
          \addplot+[domain=-3:3] {-x^2}; %
          \pgfplotsforeachungrouped \a in {-2.1, 0, 1.75} { %
            \pgfmathsetmacro{\h}{0.75}%
            \pgfmathsetmacro{\x}{\a-\h}%
            \pgfmathsetmacro{\xx}{\a+\h}%
            \pgfmathsetmacro{\y}{-\x*\x}%
            \pgfmathsetmacro{\yy}{-\xx*\xx}%
            \edef\temp{\noexpand\draw[blue, thick] (\x, \y) -- (\xx, \yy);}%
            \temp%
          }
        \end{axis}
      \end{tikzpicture}      
    \end{center}
  \end{solution}

  \begin{problem}
    Use your graphs to fill in the following blanks with either ``above'' or ``below''.
  \end{problem}

  \begin{itemize}

  \item When a curve is concave up on an interval, its secant line segments in that interval lie \fitb{0.8in}{above} the curve.

  \item When a curve is concave down on an interval, its secant line segments in that interval lie \fitb{0.8in}{below} the curve.

  \end{itemize}

\item 

  \note{This is just to point out that, if they're unsure where the change in concavity happens, they can try drawing some nearby secant lines.}  

  \begin{problem}[\nstretch{0.25}]
    Roughly where is the graph in \ref{mass-population} concave up? Roughly where is it concave down?
  \end{problem}

  \begin{solution}
    Looking at where secant line segments lie above and below the graph, it roughly appears that the
    function is concave down on the interval $(0,4)$ and concave up on the interval $(4,8)$.
  \end{solution}

\item  \label{swim-race} 

  \emph{Average speed vs.\ average velocity.} A swimmer is swimming a 100 m long race, which is one lap in a 50 m long pool. Let $s(t)$ be the swimmer's distance from the starting position $t$ seconds after the start of the race.

  \begin{enumerate}
  \item

    \note{The key point here is that $s(t)$ is showing us the swimmer's position at time $t$, not the total distance they've gone by time $t$. I will highlight this because it's the key to understanding the difference between velocity and speed.}

    \begin{problem}[0.25in]
      Which of the following is a more reasonable graph for $s(t)$?  Why? 

      \begin{center}
        \begin{tabular}{ccc}
          \begin{tikzpicture}
            \begin{axis}[gridgraph, cycle list=black, xlabel=$t$, 
              ylabel=$s$, scale=0.75, ytick={20, 40, ..., 100}, xtick={10, 
                20, 30, 40, 50}]
              \addplot+[domain=0:1]{100*(x/50)^(3/4)};
              \addplot+[domain=0:50]{100*(x/50)^(3/4)};
            \end{axis}
          \end{tikzpicture}
          & \hspace{0.5in} & 
          \begin{tikzpicture}
            \begin{axis}[ytick={10, 20, 30, 40, 50}, gridgraph, cycle
              list=black, xlabel=$t$, ylabel=$s$, scale=0.75, 
              xtick={10, 20, 30, 40, 50}] 
              \addplot+[domain=0:1]{50*(1 - (2*(x/50)^(ln(2)/ln(5/2))-1)^2)};
              \addplot+[domain=0:50]{50*(1 - (2*(x/50)^(ln(2)/ln(5/2))-1)^2)};
            \end{axis}
          \end{tikzpicture} \\
          (A) && (B) 
        \end{tabular}
      \end{center}
    \end{problem}

    \begin{solution}
      Choice (B) is the reasonable one: when the swimmer reaches the opposite end of the pool, they are 50 m from the starting point. At the end of the race, they are back at the starting point, so they are 0 m from the starting point. (Choice (A) would say that, at the end of the race, the swimmer is 100 m from his starting point.)
    \end{solution}

  \item \label{swim-race-average-rates-for-race}

    \note{Here's where I'll introduce the difference between average speed and average velocity.}

    \begin{problem}[\vfill]
      According to the graph you chose, what was the swimmer's average speed for the race? Average velocity for the race?
    \end{problem}

    \begin{solution}
      First, average speed = $\frac{\text{distance traveled}}{\text{elapsed time}}$, while average
      velocity = $\frac{\text{change in position}}{\text{elapsed time}}=\frac{\Delta s}{\Delta t}$.

      The total distance traveled by the swimmer from the start of the race to the end of the race is
      100 m, and  the race takes 50 seconds, so their average speed for the race is 
      $\frac{100 \text{ m}}{50 \text{ s}} = 2 \text{ m/s}$.

      The change in the swimmer's position from the start of the race to the end of the race
      is 0 because they end in the same place they started, so their average velocity for the race is $\frac{0 \text{ m}}{50 \text{ s}} = 0 \text{ m/s}$.
    \end{solution}

  \item
    \begin{problem}[\vfill]
      What was the swimmer's average speed over the last 50 m of the race? Average velocity?
    \end{problem}

    \begin{solution}
      It took the swimmer 30 seconds to travel the last 50 meters, so their average speed was 
      $\frac{50 \text{ m}}{30 \text{ s}} = \frac{5}{3}$ m/s.

      In the same 30 second interval, the change in the swimmer's position was $s(50)-s(20)=-50$ m, 
      so their average velocity was $\frac{\Delta s}{\Delta t}=
      \frac{s(50) - s(20)}{50 - 20} = \frac{-50 \text{ m}}{30 \text{ s}} = - \frac{5}{3} \text{ m/s}$.
      The velocity is negative because they were swimming back towards the starting position.
    \end{solution}

  \item
    \begin{problem}[\vfill]
      What was the swimmer's average speed over the first 20 seconds of the race? Average velocity?
    \end{problem}

    \begin{solution}      
      In the first 20 seconds, the swimmer swam a distance of 50 m. So, their average speed on this interval was $\frac{50 \text{ m}}{20 \text{ s}} = 2.5 \text{ m/s}$.

      The swimmer's average velocity over the first 20 seconds is $\frac{\Delta s}{\Delta t} = \frac{s(20) - s(0)}{20 - 0} = \frac{50 \text{ m}}{20 \text{ s}} = {2.5 \text{ m/s}}$.
    \end{solution}

  \item \label{swim-race-velocity-vs-speed}

    \note{I'll talk about average velocity vs.\ average speed, as well as instantaneous velocity vs.\ instantaneous speed.}

    \begin{problem}[\vfill]
      What's the difference between velocity and speed?
    \end{problem}

    \begin{solution}
      Velocity takes into account the direction of motion and can be negative. It is the
      rate of change of position with respect to time. 
      Speed is the rate of change of distance
      traveled with respect to time and is always non-negative.
      
      The average velocity and average speed can have completely different values. For example, in
      \ref{swim-race-average-rates-for-race}, the average speed for the race is 2 m/s, but the average
      velocity is 0 m/s. However, instantaneous speed is always the absolute value of
      instantaneous velocity.
    \end{solution}

  \item \label{swim-race-average-acceleration}

    \note{This is to introduce the idea of average acceleration. This can be a confusing idea for students (especially the units), so help students interpret it.}

    \begin{problem}[0.5in]
      Let $v(t)$ be the swimmer's velocity at time $t$. Interpret
      $\dfrac{v(15) - v(10)}{5}$ in words. What are its units?
    \end{problem}

    \begin{solution}
        $\frac{v(15)-v(10)}{5}=\frac{\Delta v}{\Delta t}$ is the average rate of change of
        the swimmer's velocity with respect to time between $t=10$ and $t=15$. Its units are 
        $\frac{\text{m/s}}{\text{s}}$, usually written $\text{m/s$^2$}$ and pronounced 
        ``meters per second per second'' or ``meters per second squared.''
        Another word for ``average rate of change of velocity with respect to time'' is 
        \textbf{average acceleration.}
    \end{solution}

  \end{enumerate}

  \pspace{\newpage}

\item 

  \note{This is just more practice thinking about velocity and speed. The goal is to have students ``narrate'' the situation. For example:
    \begin{itemize}
    \item The height starts at 1.5, goes up, then goes down.
    \item The velocity starts out positive (the ball is going up), then switches to negative (the ball is going down). This could be (C) or (F), but (C) says that the ball goes up at a constant speed, while we know the ball should really slow down as it's going up.
    \item The ball should slow down as it goes up and speed up as it goes down, with a speed of 0 when it's at the very top. Alternatively, when the ball is going up, its speed should be exactly the same as its velocity.
    \end{itemize}
  }

  From your height of 1.5 m, you throw a ball straight up into the air. Pick the graph (choices A - G below) that could be the ball's\ldots
  \begin{enumerate}
  \item
    \begin{problem}[\vfill]
      height as a function of time (for the first few moments after you toss the ball)
    \end{problem}

    \begin{solution}
      The height of the ball should start at 1.5 m, increase, then decrease. This matches
      graph (D).
    \end{solution}

  \item
    \begin{problem}[\vfill]
      velocity as a function of time
    \end{problem}

    \begin{solution}
      The velocity will initially be positive as the ball moves upwards, decrease until it reaches
      zero when the ball reaches its maximum height, then become negative as the ball falls towards
      the ground.
      This matches graph (F).
    \end{solution}

  \item
    \begin{problem}[\vfill]
      speed as a function of time
    \end{problem}

    \begin{solution}
      The speed will initially be positive as the ball moves upwards, then decrease until it
      reaches zero when the ball reaches its maximum height, then increase again as the ball
      falls towards the ground. This matches graph (E). Instantaneous speed is the absolute value
      of instantaneous velocity.
    \end{solution}

  \end{enumerate}

  {
    \providecommand{\choice}{}
    \renewcommand{\choice}[2]{
      \begin{tikzpicture}
        \begin{axis}[xtick=\empty, ytick=\empty, width=1.65in, xlabel=$t$, height=1.65in, xmax=2.2, #2]
          \addplot+[domain=0:2] {#1};
        \end{axis}
      \end{tikzpicture}
    }

    \begin{imagetable}{4}{\Alph}
      \choice{2-2*x+x^2}{ymin=0} & \choice{1-abs(x-1)}{ymax=2} &
      \begin{tikzpicture}
        \begin{axis}[xtick=\empty, ytick=\empty, width=1.65in, xlabel=$t$, ymax=2, ymin=-2, height=1.65in, xmax=2.2]
          \addplot+[domain=0:1] {1}; %
          \addplot+[domain=1:2] {-1}; %
          \addplot+[open] coordinates { (1, 1) (1, -1) }; %
        \end{axis}
      \end{tikzpicture}
      & \choice{1 + 2*x - x^2}{ymin=0} \\
      \begin{tikzpicture}
        \begin{axis}[xtick=\empty, ytick=\empty, width=1.65in, xlabel=$t$, height=1.65in, xmax=2.2]
          \addplot+[domain=0:2, opacity=0] {2-2*x}; % to get the same scale as the next graph
          \addplot+[domain=0:2] { abs(2-2*x)}; 
        \end{axis}
      \end{tikzpicture}
      & \choice{2-2*x}{} & \choice{x}{}
    \end{imagetable}
  }

  \pspace{\newpage}

\item 

  In this problem, we'll look at the function $f(x) = x^2$.  Calculate the
  average rate of change of $f$ on the following intervals.
  \begin{enumerate}
  \item
    \begin{problem}[\vfill]
      $[1, 3]$
    \end{problem}

    \begin{solution}
      $\dfrac{f(3)-f(1)}{3-1}=\dfrac{3^2-1}{3-1}=\dfrac{9-1}{3-1}=\dfrac{8}{2}=\boxed{4}$
    \end{solution}

  \item
    \begin{problem}[\vfill]
      $[1, 1.5]$
    \end{problem}

    \begin{solution}
      $\dfrac{f(1.5)-f(1)}{1.5-1}=\dfrac{1.5^2-1}{1.5-1} = \dfrac{2.25-1}{1.5-1}= 
      \dfrac{1.25}{0.5}=\boxed{2.5}$
    \end{solution}

  \item

    \note{Practice with the algebra we'll do when finding derivatives using the definition of the derivative.}

    \begin{problem}[\stretch{2}]
      $[1, 1 + h]$.  (From your work in (a) and (b), you should already
      know the answer for certain values of $h$; which ones?  Does your
      answer here agree with your answers to (a) and (b)?)
    \end{problem}

    \begin{solution}
      $\dfrac{f(1+h)-f(1)}{1+h-1}=\dfrac{(1+h)^2-1}{1+h-1}=\dfrac{1+2h+h^2-1}{1+h-1}
      = \dfrac{2h+h^2}{h} =\boxed{2+h, h \neq 0}$
    \end{solution}

  \item
    \begin{problem}[\stretch{2}]
      Sketch the graph of $f(x) = x^2$, and draw secant lines whose slopes
      correspond to the average rates of change in (a) and (b).  Which of
      the average rates of change is larger?
    \end{problem}

    \begin{solution}
      The average rate of change on $[1, 3]$ is larger than the average rate of change
      on $[1, 1.5]$.
      
      \begin{tikzpicture}
        \begin{axis}[xmin=0, xmax=4, xtick={1,1.5,3}, ytick=\empty]
          \addplot[domain=0:4]{x^2};
          \addplot[closed] coordinates {(1,1) (1.5,2.25) (3,9)};
          \addplot[red, domain=0:4]{2.5*(x-1)+1};
          \addplot[red, domain=0:4]{4*(x-1)+1};
        \end{axis}
      \end{tikzpicture}
    \end{solution}

  \end{enumerate}

\item  

  \note{We \emph{could} calculate, but if we think about these rates
    as slopes, then it's obvious they're all the same.}

  \begin{problem}[\stretch{2}]
    Suppose $f(x)$ is a linear function defined on $(-\infty, 
    \infty)$. How does the average rate of change of $f$ on $[0, 5]$
    compare to the average rate of change of $f$ on $[0, 500]$? What about
    the average rate of change of $f$ on $[-50, 2]$?
  \end{problem}

  \begin{solution}
    They are all equal to the slope of the line!
  \end{solution} 

\end{enumerate}
\end{document}
